\subsection{Burnol's local jets in the Weil completion}
\label{subsec:fd-burnol-weil}

Burnol's local vectors carry a full jet Gram matrix in physical
\(L^2\). The following computation describes their image under the
Weil sampling map and the term that remains when these vectors are
localized to a finite window.

Write \(\mathcal K=L^2((0,\infty),dt)\), and let
\[
 (\mathcal J f)(x)=e^{x/2}f(e^x).
\]
This is unitary from \(\mathcal K\) to \(L^2(\mathbb R,dx)\).
With the Fourier convention of this paper,
\(\widehat{\mathcal J f}(\nu)=\operatorname{Mel}f(1/2-i\nu)\)
where either side is defined. Burnol's unitary \(V\) on
\(\mathcal K\) has Mellin multiplier
\[
 v(s)=\left(\frac{s}{1-s}\right)^3
                \frac{\zeta(1-s)}{\zeta(s)}
 =\left(\frac{s}{1-s}\right)^3
       2(2\pi)^{-s}\Gamma(s)\cos\frac{\pi s}{2}.
\]
The second expression removes the apparent singularities at zeta
zeros, and \(|v(1/2+it)|=1\).

Fix a signed ordinate \(\gamma\), set \(\rho=1/2-i\gamma\),
and let \(0<\lambda\le e^{-1}\), \(R=\log(1/\lambda)\).
Let \(Q_\lambda\) be multiplication by
\(\mathbf1_{[\lambda,\infty)}\). Burnol's vectors have the
form
\[
 Y_{\rho,k}^{\lambda}=V^{-1}Q_\lambda\Phi_{\rho,k},
 \qquad
 X_k^\lambda=R^{-k-1/2}Y_{\rho,k}^{\lambda},
 \qquad k\ge0.
\]
Here \(\Phi_{\rho,k}\) is the boundary value of
\(V((\log(1/t))^k t^{-w}\mathbf1_{(0,1)}(t))\) as
\(w\to\rho\) from the left. The existence and physical
estimates are Burnol's Theorem~4.2 and Theorem~4.11
\cite[pp.~9--13]{burnol2002}. In his multiplicity application one
uses \(0\le k<m_\gamma\), the range for which the vectors have
the required orthogonality to the approximation space.

We use the following explicit consequences of that construction:
\begin{align}
 \operatorname{Mel}\Phi_{\rho,k}(s)
  &=\frac{k!v(s)}{(s-\rho)^{k+1}}
               &&(1/2<\Re s<1),\label{eq:fd-burnol-mellin}\\
 \Phi_{\rho,k}(t)
  &=\partial_\rho^k\bigl(v(\rho)t^{-\rho}\bigr)
       +r_{\rho,k}(t)
               &&(0<t\le1),\label{eq:fd-burnol-small}\\
 |r_{\rho,k}(t)|+|t r_{\rho,k}'(t)|&\le C_{\rho,k}t^2,
               &&(0<t\le1),\label{eq:fd-burnol-rem}\\
 \Phi_{\rho,k}(t)&=O_{\rho,k}((1+\log t)^2/t)
               &&(t\ge1).\label{eq:fd-burnol-large}
\end{align}
For clarity, the derivative bound in
\eqref{eq:fd-burnol-rem} follows from the convergent series
preceding Burnol's Theorem~4.11, followed by the two Hardy averages
in \(V=(1-M)^2U\). The displayed source series is stated for
\(k\ge1\). For \(k=0\), integration by parts in Burnol's
defining integral gives the boundary term
\(\sin(2\pi t)/(\pi t)\). Its constant term \(+2\) cancels
the constant \(-2\) from extending the remaining integral to
zero, so the same series starting at \(n=1\) holds. The
resulting formula for every \(k\ge0\) is
\[
 r_{\rho,k}(t)=2(-1)^k k!\sum_{n\ge1}
 \frac{(-1)^n(2\pi t)^{2n}}{(2n+1)!}
 \frac{2n}{(\rho+2n)^{k+1}}
 \left(\frac{2n}{2n+1}\right)^2.
\]
The series and its \(t\partial_t\) derivative converge absolutely
on \([0,1]\) and start with \(t^2\). The Mellin identity
\eqref{eq:fd-burnol-mellin} first follows from the multiplier
identity for \(\Re w<1/2\); the physical bounds permit passage
to \(w=\rho\) on each closed substrip of \(1/2<\Re s<1\).

\begin{proposition}[The leading Weil image of the local jets]
\label{prop:fd-burnol-weil}
The canonical Fourier transforms of \(\mathcal JX_k^\lambda\)
define vectors in the Weil completion. For each fixed \(\gamma\)
and \(k\),
\begin{equation}\label{eq:fd-burnol-coordinate}
 \left\|\mathcal JX_k^\lambda
       -\frac{\sqrt R}{k+1}\mathbf1_{\{\gamma\}}\right\|_q
       =O_{\gamma,k}(R^{-1/2}).
\end{equation}
The coordinate vector is in \(\ell^2(\Gamma,m)\), so its squared
norm is \(m_\gamma\). In particular,
\begin{equation}\label{eq:fd-burnol-weil-gram}
 R^{-1}q(\mathcal JX_j^\lambda,\mathcal JX_k^\lambda)
 \longrightarrow\frac{m_\gamma}{(j+1)(k+1)}.
\end{equation}
For every fixed finite coefficient vector \(a=(a_0,\ldots,a_{r-1})\)
with \(\sum_k a_k/(k+1)=0\),
\[
 q\left(\sum_{k<r}a_k\mathcal JX_k^\lambda\right)
       =O_{\gamma,r}(R^{-1})\|a\|_2^2.
\]
Its physical squared norm tends to
\(\sum_{j,k<r}a_j\overline{a_k}/(j+k+1)\), which is positive
for \(a\ne0\).
\end{proposition}

\begin{proof}
Put \(s_\nu=1/2-i\nu\), and define
\[
 E_\lambda(s)=\int_0^\lambda r_{\rho,k}(t)t^{s-1}\,dt.
\]
Integration by parts, using \eqref{eq:fd-burnol-rem}, gives on
the critical line
\begin{equation}\label{eq:fd-burnol-cutoff-error}
 |E_\lambda(s_\nu)|\le
                 C_{\rho,k}\lambda^{5/2}/|s_\nu|.
\end{equation}
Indeed the boundary term is \(\lambda^s r(\lambda)/s\), and
the remaining integral is \(-s^{-1}\int_0^\lambda r'(t)t^sdt\).
Both have the displayed bound.

Subtracting the lower tail from
\eqref{eq:fd-burnol-mellin} gives the canonical transform
\begin{align}
 F_{R,k}(\nu)
 &:={\widehat{\mathcal JY_{\rho,k}^{\lambda}}}(\nu)\notag\\
 &=\frac1{v(s_\nu)}\left\{
   \frac{k!v(s_\nu)}{(s_\nu-\rho)^{k+1}}
   -\partial_\rho^k
       \left(\frac{v(\rho)\lambda^{s_\nu-\rho}}
                         {s_\nu-\rho}\right)
   -E_\lambda(s_\nu)\right\}.
 \label{eq:fd-burnol-transform}
\end{align}
At \(\nu=\gamma\), the apparent singularity is removable.
One may obtain the value directly by integrating
\(Q_\lambda\Phi_{\rho,k}\). Equations
\eqref{eq:fd-burnol-small}--\eqref{eq:fd-burnol-large} give
\[
 \int_\lambda^\infty\Phi_{\rho,k}(t)t^{\rho-1}\,dt
 =\sum_{j=0}^k\binom kj v^{(j)}(\rho)
        \frac{R^{k-j+1}}{k-j+1}+O_{\rho,k}(1).
\]
Division by \(v(\rho)\) therefore yields
\begin{equation}\label{eq:fd-burnol-target}
 F_{R,k}(\gamma)=\frac{R^{k+1}}{k+1}+O_{\gamma,k}(R^k).
\end{equation}

For a distinct ordinate \(\nu\), the difference
\(|s_\nu-\rho|=|\nu-\gamma|\) is bounded below by a positive
constant depending on \(\gamma\). The derivatives of
\(\lambda^{s-\rho}/(s-\rho)\) are finite sums of
\[
 \lambda^{s-\rho}R^{\ell-j}j!(s-\rho)^{-j-1},
 \qquad 0\le j\le\ell\le k.
\]
On the critical line the exponential factor has modulus one.
Since \(|v(s_\nu)|=1\), equations
\eqref{eq:fd-burnol-cutoff-error} and
\eqref{eq:fd-burnol-transform} imply
\begin{equation}\label{eq:fd-burnol-other}
 |F_{R,k}(\nu)|\le
    C_{\gamma,k}\left(
       \frac{R^k}{|\nu-\gamma|}
           +\frac{\lambda^{5/2}}{|s_\nu|}\right).
\end{equation}
Both squared majorants are summable against the zero weights:
\[
 \sum_{\nu\ne\gamma}\frac{m_\nu}{(\nu-\gamma)^2}<\infty,
 \qquad
 \sum_\nu\frac{m_\nu}{|s_\nu|^2}<\infty.
\]
This follows from isolation of the fixed zero and the counted
unit-interval estimate. Multiplying
\eqref{eq:fd-burnol-target} and
\eqref{eq:fd-burnol-other} by \(R^{-k-1/2}\) proves
\eqref{eq:fd-burnol-coordinate}, provided the canonical transform
is interpreted in the sampling completion. We justify that
interpretation next.

For each fixed \(\lambda,\rho,k\),
\eqref{eq:fd-burnol-transform} is a continuous function of \(\nu\)
with a continuous derivative, including at the removable point.
The gamma expression for \(v\) gives
\[
 |v'(1/2-i\nu)|=O(\log(2+|\nu|)).
\]
For example this follows by logarithmic differentiation: the
digamma series, split at \(n=\lceil|\nu|\rceil+1\), bounds
the digamma term by \(O(\log(2+|\nu|))\), and the remaining
rational and tangent terms are bounded outside a fixed interval.
The derivative in \(\nu\) of the cutoff error satisfies the
same \((1+|\nu|)^{-1}\) bound with a fixed extra factor
\(R+1\), by the preceding integration by parts applied to
\(r(t)\log t\). Consequently, as \(|\nu|\to\infty\),
\[
 |F_{R,k}(\nu)|+|F_{R,k}'(\nu)|
       =O_{\lambda,\gamma,k}
          \left(\frac{\log(2+|\nu|)}{1+|\nu|}\right).
\]
Thus this transform belongs to
\[
 \mathcal E=\left\{F:\int_{\mathbb R}\log(3+|t|)
                    (|F(t)|^2+|F'(t)|^2)dt<\infty\right\}.
\]
The unit-interval sampling inequality is continuous on
\(\mathcal E\). Smooth frequency cutoffs and mollification show
that Schwartz functions are dense in this weighted Sobolev space.
Their inverse Fourier transforms are Schwartz; cutting these off
smoothly on the physical line then converges in the Schwartz
topology and hence in \(\mathcal E\). The resulting compact
smooth tests therefore converge to the displayed sampling vector
in Weil norm. This proves the stated completion-domain assertion.

Equation \eqref{eq:fd-burnol-weil-gram} and the coefficient
annihilation estimate now follow from
\eqref{eq:fd-burnol-coordinate} and the triangle inequality in
the completed Hilbert space.

For the physical Gram, unitarity of \(V\) and \(\mathcal J\)
reduces the calculation to the functions
\(Q_\lambda\Phi_{\rho,k}\). Their leading terms on
\((\lambda,1)\) are
\(v(\rho)t^{-\rho}(\log(1/t))^k\). The leading cross integral
is exactly
\[
 |v(\rho)|^2\int_\lambda^1
       (\log(1/t))^{j+k}\frac{dt}{t}
 =\frac{R^{j+k+1}}{j+k+1}.
\]
The lower logarithmic powers contribute
\(O(R^{j+k})\), and the remainder and upper tail contribute
lower order terms by
\eqref{eq:fd-burnol-rem}--\eqref{eq:fd-burnol-large}.
Normalization gives the Hilbert Gram limit, which is Burnol's
Theorem~5.3 in this notation \cite[p.~14]{burnol2002}.
Its positive definiteness follows directly from
\[
 \sum_{j,k<r}\frac{a_j\overline{a_k}}{j+k+1}
 =\int_0^1\left|\sum_{k<r}a_k t^k\right|^2dt>0
                         \quad(a\ne0).
\]
\end{proof}

This calculation gives an exact norm comparison: the normalized
physical jet Gram has full rank, while its leading Weil image is
the single value coordinate at the chosen ordinate. The calculation
uses no upper bound for a reciprocal zeta derivative. The
multiplicity restriction \(k<m_\gamma\) enters Burnol's
orthogonality theorem; the sampling estimates themselves apply to
every fixed \(k\).

\begin{proposition}[The physical window occupied by the jets]
\label{prop:fd-burnol-physical-window}
For each fixed \(\gamma\) and \(k\), put
\[
 p_{R,k}(x)=R^{-k-1/2}(-x)^k e^{i\gamma x}
                            \mathbf1_{[-R,0]}(x).
\]
Then
\begin{equation}\label{eq:fd-burnol-physical-window}
 \|\mathcal JX_k^{e^{-R}}-p_{R,k}\|_{L^2(\mathbb R)}
                  =O_{\gamma,k}(R^{-1/2}).
\end{equation}
For a fixed finite coefficient vector \(a\), set
\(P(t)=\sum_k a_k t^k\). Uniformly for \(0\le L\le R\),
\[
 \left\|\mathbf1_{[-L,L]}
             \sum_k a_k\mathcal JX_k^{e^{-R}}\right\|_2^2
       =\int_0^{L/R}|P(t)|^2dt+O_{\gamma,a}(R^{-1/2}).
\]
In particular, a window with \(L=o(R)\) retains vanishing
physical mass, although the full physical squared norm tends to
\(\int_0^1|P(t)|^2dt\).
\end{proposition}

\begin{proof}
The Fourier transform of the unnormalized polynomial indicator is
\[
 P_{R,k}(\nu)=\frac{k!}{(s_\nu-\rho)^{k+1}}
        -\partial_\rho^k\left(
                 \frac{\lambda^{s_\nu-\rho}}{s_\nu-\rho}\right).
\]
This follows by integrating \((-x)^k e^{i\gamma x}\) on
\([-R,0]\), with the removable value at \(\nu=\gamma\).
Subtract this expression from \eqref{eq:fd-burnol-transform}.
The result is
\[
 F_{R,k}(\nu)-P_{R,k}(\nu)
 =\frac1{v(s_\nu)}\partial_\rho^k
       \left\{\lambda^{s_\nu-\rho}
          \frac{v(s_\nu)-v(\rho)}{s_\nu-\rho}\right\}
       -\frac{E_\lambda(s_\nu)}{v(s_\nu)}.
\]
For each fixed \(j\), the divided difference
\[
 Q_j(\nu)=\partial_\rho^j
                 \frac{v(s_\nu)-v(\rho)}{s_\nu-\rho}
\]
has a removable singularity at \(\nu=\gamma\) and is bounded
on compact intervals. As \(|\nu|\to\infty\), the fact that
\(|v(s_\nu)|=1\), together with the fixed derivatives of
\(v\) at \(\rho\), gives \(Q_j(\nu)=O_{\gamma,j}(|\nu|^{-1})\).
Consequently \(Q_j\in L^2(\mathbb R,d\nu)\).
Differentiating the exponential factor yields powers of \(R\),
so the last display and \eqref{eq:fd-burnol-cutoff-error} imply
\[
 \|F_{R,k}-P_{R,k}\|_{L^2(d\nu)}
      \le \sum_{j=0}^k\binom kj R^{k-j}\|Q_j\|_2
                       +O_{\gamma,k}(e^{-5R/2})
      =O_{\gamma,k}(R^k).
\]
Divide by \(R^{k+1/2}\) and apply Fourier Plancherel to obtain
\eqref{eq:fd-burnol-physical-window}. A finite linear combination
has the same error order. Multiplication by
\(\mathbf1_{[-L,L]}\) is a contraction in \(L^2\), while direct
substitution \(t=-x/R\) gives exactly
\[
 \left\|\mathbf1_{[-L,L]}\sum_k a_kp_{R,k}\right\|_2^2
                      =\int_0^{L/R}|P(t)|^2dt.
\]
The squared-norm difference is bounded by the norm error times
the sum of the two norms, which is uniformly bounded. This proves
the assertion uniformly in \(L\).
\end{proof}

The remaining finite-window term is explicit. For a smooth physical
cutoff \(\chi_L\) and
\(F_R=\sum_{k<r}a_k\mathcal JX_k^\lambda\), with
\(\sum a_k/(k+1)=0\), the triangle inequality gives
\[
 q(\chi_LF_R)^{1/2}
 \le C_{\gamma,r}R^{-1/2}\|a\|_2
             +q((1-\chi_L)F_R)^{1/2}.
\]
For each fixed \(R\), smooth cutoffs approaching one converge in
the sampling domain after smoothing if necessary. One direct
justification uses their approximate-identity convolution in
frequency and the moderate weight \(\log(3+|t|)\); weighted
Young's inequality gives uniform boundedness, and density gives
strong convergence in \(\mathcal E\), also for the first
derivative. The following calculation gives the cutoff norm for a fixed
polynomial family at a fixed ordinate. A bound uniform in the
height and in the number of jet orders, at the scales needed for
the all-height criterion, remains unproved.

\subsection{The Weil norm of a coupled physical cutoff}
\label{subsec:burnol-coupled-cutoff}
The local vectors of
Proposition~\ref{prop:fd-burnol-weil} admit a quantitative
cutoff calculation for each fixed finite polynomial family.
Both endpoints of the physical polynomial enter the
uncut transform; the endpoint at zero determines the leading
sample after a shorter physical cutoff.

Fix $\gamma\in\Gamma$, put $\rho=1/2-i\gamma$, and retain
Burnol's multiplier
\[
 v(s)=\left(\frac{s}{1-s}\right)^3
       2(2\pi)^{-s}\Gamma(s)\cos\frac{\pi s}{2}.
\]
For a polynomial $P(t)=\sum_{k=0}^K b_kt^k$, define
\[
 H_R=\sum_{k=0}^K b_k\mathcal JX_k^{e^{-R}},
 \qquad
 p_R(x)=R^{-1/2}P(-x/R)e^{i\gamma x}
                    {\bf1}_{[-R,0]}(x).
\]
All constants below may depend on this fixed ordinate and
polynomial. Choose a real nonnegative
$\chi\in C_c^\infty((-1,1))$ equal to one near zero, and set
$\chi_L(x)=\chi(x/L)$.

\begin{proposition}\label{prop:burnol-coupled-cutoff}
Suppose $\int_0^1P(t)\,dt=0$, that $P$ has exact order $M\ge0$
at zero, and that it vanishes to order at least $M$ at one.
Write
\[
 a=\frac{P^{(M)}(0)}{M!}\ne0,\quad
 A_{\chi,M}=\int_0^1\chi(-t)t^Mdt,\quad
 B_{\chi,M}=\int_0^1\chi(-t)^2t^{2M}dt.
\]
Then $q(H_R)=O_{\gamma,P}(R^{-2M-1})$. For every
$L=L(R)\to\infty$ with $L/R\to0$,
\begin{align}
 q(\chi_LH_R)
 &\sim q((1-\chi_L)H_R)
 \sim m_\gamma|a|^2A_{\chi,M}^2
                  \frac{L^{2M+2}}{R^{2M+1}},
                 \label{eq:burnol-cutoff-q}\\
 \|\chi_LH_R\|_2^2
 &\sim |a|^2B_{\chi,M}
                  \frac{L^{2M+1}}{R^{2M+1}}.
                 \label{eq:burnol-cutoff-physical}
\end{align}
In particular,
\[
 \frac{q(\chi_LH_R)}{\|\chi_LH_R\|_2^2}
 \sim m_\gamma\frac{A_{\chi,M}^2}{B_{\chi,M}}\,L.
\]
\end{proposition}

\begin{proof}
We first record the endpoint cancellations in the
multiplier formula. Put $s_\nu=1/2-i\nu$ and
$d_\nu=s_\nu-\rho=-i(\nu-\gamma)$. Summing
\eqref{eq:fd-burnol-transform} against $b_kR^{-k-1/2}$
and applying Leibniz's rule gives, for $\nu\ne\gamma$,
\begin{align}
 \widehat H_R(\nu)
 =R^{-1/2}\bigg\{&
 \sum_{j=0}^K\frac{P^{(j)}(0)}{R^jd_\nu^{j+1}}\notag\\
 &-\frac{e^{-Rd_\nu}}{v(s_\nu)}
 \sum_{j=0}^K\frac{P^{(j)}(1)}{j!R^j}
       \partial_\rho^j\left(\frac{v(\rho)}{d_\nu}\right)
 \bigg\}+\mathcal R_R(\nu).
 \label{eq:burnol-endpoint-transform}
\end{align}
Indeed $\sum_{k\ge j}b_k\binom{k}{j}=P^{(j)}(1)/j!$,
while $b_jj!=P^{(j)}(0)$. The lower-tail remainder bound
\eqref{eq:fd-burnol-cutoff-error} gives
\begin{equation}\label{eq:burnol-poly-rem}
 |\mathcal R_R(\nu)|
 \le C_{\gamma,P}R^{-M-1/2}
           e^{-5R/2}/|s_\nu|.
\end{equation}
Here $b_k=0$ for $k<M$. The two endpoint sums in
\eqref{eq:burnol-endpoint-transform} start at $j=M$.
Isolation of the fixed ordinate and $|v(s_\nu)|=1$ therefore
give
\[
 |\widehat H_R(\nu)|
 \le C_{\gamma,P}R^{-M-1/2}|\nu-\gamma|^{-1}
      +|\mathcal R_R(\nu)|\quad(\nu\in\Gamma,\ \nu\ne\gamma).
\]
The squared majorants are summable against $m_\nu$.

At the target coordinate the direct integral calculation
used for \eqref{eq:fd-burnol-target} instead gives
\[
 \widehat H_R(\gamma)=
 \sqrt R\sum_{j=0}^K
 \frac{v^{(j)}(\rho)}{v(\rho)j!R^j}
             \int_0^1P^{(j)}(t)\,dt
          +O_{\gamma,P}(R^{-M-1/2}).
\]
The $j=0$ term vanishes by the integral constraint. For
$1\le j\le M$, the integral equals
$P^{(j-1)}(1)-P^{(j-1)}(0)=0$. Every remaining term is
$O(R^{-M-1/2})$. This proves $q(H_R)=O(R^{-2M-1})$.
It also identifies the two leading off-target endpoint terms:
for each fixed $\nu\ne\gamma$,
\[
 \widehat H_R(\nu)=
 \frac{R^{-1/2}}{-i(\nu-\gamma)}
 \left\{P(0)-P(1)\frac{v(\rho)}{v(s_\nu)}
                     e^{iR(\nu-\gamma)}\right\}
 +O_{\gamma,P,\nu}(R^{-3/2}).
\]

We next control multiplication by $\chi_L$ in the sampling
norm. Set $w(t)=\log(3+|t|)$ and
$\|F\|_w^2=\int w(t)|F(t)|^2dt$.
The counted unit-interval zero bound implies, for
$0<\epsilon\le1$,
\begin{equation}\label{eq:burnol-short-sampling}
 \sum_{\nu\in\Gamma}m_\nu|F(\nu)|^2
 \le C\left(\epsilon^{-1}\|F\|_w^2
                         +\epsilon\|F'\|_w^2\right).
\end{equation}
To prove it, apply the elementary one-dimensional Sobolev
inequality at $\nu$ on $[\nu-\epsilon,\nu+\epsilon]$.
At a point $t$, at most $C\log(3+|t|)$ counted ordinates
have intervals containing it. Summation and integration
give \eqref{eq:burnol-short-sampling}.

If $e$ has Fourier transform in $L^2(w\,dt)$, multiplication
by $\chi_L$ is convolution with
$k_L(t)=L\widehat\chi(Lt)/(2\pi)$ in frequency.
The inequality
$w(t+u)\le Cw(t)\log(3+|u|)$ and Minkowski's inequality give
\[
 \|k_L*F\|_w\le C_\chi\|F\|_w,\qquad
 \|k_L'*F\|_w\le C_\chi L\|F\|_w
 \qquad(L\ge1).
\]
Both follow by substituting $uL$ in the weighted $L^1$
norms of the Schwartz kernels. Taking $\epsilon=1/L$ in
\eqref{eq:burnol-short-sampling} proves the useful bound
\begin{equation}\label{eq:burnol-cutoff-sampling}
 q(\chi_Le)\le C_\chi L\|\widehat e\|_w^2.
\end{equation}
This explicitly uses weighted Fourier control; a physical
$L^2$ bound alone has not been substituted for a Weil bound.

The exact difference from the polynomial indicator is
\begin{equation}\label{eq:burnol-poly-error}
 \widehat H_R(\nu)-\widehat p_R(\nu)
 =\frac{R^{-1/2}e^{-Rd_\nu}}{v(s_\nu)}
 \sum_{j=M}^K\frac{P^{(j)}(1)}{j!R^j}\,Q_j(\nu)
       +\mathcal R_R(\nu),
\end{equation}
where
\[
 Q_j(\nu)=\partial_\rho^j
             \frac{v(s_\nu)-v(\rho)}{s_\nu-\rho}.
\]
Each singularity at $\gamma$ is removable, and
$Q_j(\nu)=O_{\gamma,j}(|\nu|^{-1})$ at infinity.
Consequently $Q_j\in L^2(w\,d\nu)$. On the real line the
two factors $e^{-Rd_\nu}$ and $v(s_\nu)^{-1}$ have absolute
value one. Thus
\[
 \|\widehat H_R-\widehat p_R\|_w
       \le C_{\gamma,P}R^{-M-1/2}.
\]
Equations \eqref{eq:burnol-cutoff-sampling} and Fourier
Plancherel now give
\begin{equation}\label{eq:burnol-two-cutoff-errors}
 q\bigl(\chi_L(H_R-p_R)\bigr)
       \le C_{\gamma,P,\chi}L R^{-2M-1},\qquad
 \|\chi_L(H_R-p_R)\|_2
       \le C_{\gamma,P}R^{-M-1/2}.
\end{equation}
These statements also place the compactly supported
functions in the completed form domain. For example,
physical mollification supported within the fixed support
margin converges in the weighted Fourier Sobolev norm
appearing in \eqref{eq:burnol-short-sampling}; it gives
compact smooth approximants in the same window.

It remains to compute the compact polynomial itself.
Write $g_{R,L}=\chi_Lp_R$. At the target,
\begin{align*}
 \widehat g_{R,L}(\gamma)
 &=R^{-1/2}L\int_0^1\chi(-t)P(Lt/R)\,dt\\
 &=aA_{\chi,M}\frac{L^{M+1}}{R^{M+1/2}}
                      \left(1+O_P(L/R)\right).
\end{align*}
At a distinct ordinate, integration by parts $M+1$ times
gives
\[
 |\widehat g_{R,L}(\nu)|
 \le C_{P,\chi}R^{-M-1/2}|\nu-\gamma|^{-M-1}
       \qquad(1\le L\le R/2).
\]
For completeness, all left endpoint terms vanish because
$\chi$ has compact support in $(-1,1)$. At zero the
derivatives of order less than $M$ vanish, the derivative
of order $M$ is $(-1)^MP^{(M)}(0)/R^M$, and the
$L^1$ norm of the derivative of order $M+1$ of
$\chi(x/L)P(-x/R)$ on $[-L,0]$ is $O_{P,\chi}(R^{-M})$.
The last bound follows term by term from Leibniz's rule,
using $P^{(j)}(t)=O_P(t^{M-j})$ for $j\le M$ and $L/R\le1/2$.
Therefore the off-target squared samples sum to
$O_{\gamma,P,\chi}(R^{-2M-1})$, whereas the target contributes
the leading expression in \eqref{eq:burnol-cutoff-q}.
The first error norm in \eqref{eq:burnol-two-cutoff-errors}
is smaller than this leading norm by
$O(L^{-M-1/2})$. This proves the asymptotic for
$q(\chi_LH_R)$.
The same leading sample has the exact global jet image
given by \eqref{eq:hb-evaluation-jet}:
\[
 (\mathcal U(\chi_LH_R))^{(m_\gamma-1)}(\gamma)
 =\frac{iF^{(m_\gamma)}(\gamma)}{\sqrt\pi}
                      \widehat{\chi_LH_R}(\gamma).
\]
Thus the coordinate responsible for the cutoff norm is
identified with its corresponding derivative evaluation,
with the multiplicity and normalization retained.

Direct substitution also gives
\[
 \|g_{R,L}\|_2^2
 =R^{-1}L\int_0^1\chi(-t)^2|P(Lt/R)|^2dt
 \sim |a|^2B_{\chi,M}\frac{L^{2M+1}}{R^{2M+1}}.
\]
The second error bound in \eqref{eq:burnol-two-cutoff-errors}
is smaller than its norm by $O(L^{-M-1/2})$, proving
\eqref{eq:burnol-cutoff-physical}. Finally
$\|H_R\|_q=O(R^{-M-1/2})$ is smaller than
$\|\chi_LH_R\|_q$ by $O(L^{-M-1})$.
The reverse triangle inequality proves that
$q((1-\chi_L)H_R)\sim q(\chi_LH_R)$.
\end{proof}

\paragraph{An exponential estimate for the local multiplier error.}
The multiplier gives a stronger version of the first
bound in \eqref{eq:burnol-two-cutoff-errors}. With the same
fixed data, for $1\le L\le R/2$,
\begin{equation}\label{eq:burnol-exp-local-error}
 \|\chi_L(H_R-p_R)\|_q
 \le C_{\gamma,P,\chi}R^{-M-1/2}
 \left\{L^{3/2}e^{-2(R-L)}
                    +\sqrt L\,e^{-5R/2}\right\}.
\end{equation}
Here and below $\|f\|_q=q(f)^{1/2}$.
We give the multiplier justification to keep the direction
of the physical translation explicit.

Let $\kappa_j$ be the inverse Fourier transform of
$Q_j(\nu)/v(s_\nu)$. For $x>0$,
\begin{equation}\label{eq:burnol-positive-kernel-decay}
 |\kappa_j(x)|+|\kappa_j'(x)|
                   \le C_{\gamma,j}e^{-2x}.
\end{equation}
Indeed,
\[
 \frac{Q_j(\nu)}{v(s_\nu)}
 =\frac{j!}{d_\nu^{j+1}}
  -\frac1{v(s_\nu)}
        \partial_\rho^j\left(\frac{v(\rho)}{d_\nu}\right).
\]
The apparent real pole cancels in the whole expression.
The reciprocal multiplier is holomorphic in
$0\le\Im\nu<5/2$: its apparent issue at $s=1$ is a zero
of order two, and its first pole above the critical line
is at $s=3$, corresponding to $\nu=5i/2$.
The gamma formula and Stirling's formula give
$|v(5/2-it)^{-1}|\le C(1+|t|)^{-2}$.
Shift the inverse Fourier contour to $\nu=t+2i$.
The vertical edge terms vanish because the whole integrand
before the Fourier exponential is $O(1/|t|)$ uniformly
across this fixed strip. On the shifted line the inverse
transform of $j!/d_\nu^{j+1}$ is zero for $x>0$.
One can check this without a contour closure from
\[
 \frac{j!}{(2-i(t-\gamma))^{j+1}}
     =\int_0^\infty u^je^{-2u}e^{i(t-\gamma)u}\,du:
\]
its inverse transform is supported on $x\le0$.
The remaining shifted-line integral has integrand
$O((1+|t|)^{-3})$, and its derivative in $x$ has an
integrable $O((1+|t|)^{-2})$ majorant. Its factor $e^{-2x}$
proves \eqref{eq:burnol-positive-kernel-decay}, with a
continuous first derivative for $x>0$.

The nonremainder part of \eqref{eq:burnol-poly-error}
therefore has inverse transform
\[
 R^{-1/2}e^{-iR\gamma}
       \sum_{j=M}^K\frac{P^{(j)}(1)}{j!R^j}
                                     \kappa_j(x+R).
\]
On $|x|\le L$ its value and first derivative are bounded
by $CR^{-M-1/2}e^{-2(R-L)}$.
After multiplication by $\chi_L$ its physical $H^1$ norm
is at most $C\sqrt L R^{-M-1/2}e^{-2(R-L)}$.
The compact-support sampling estimate
$\|f\|_q\le CL\|f\|_{H^1}$ follows from
\eqref{eq:burnol-short-sampling} with $\epsilon=1$,
Fourier Plancherel, and
$w(\nu)\le C(1+\nu^2)$ applied to $f$ and $xf$.
This proves the first term in
\eqref{eq:burnol-exp-local-error}. The remainder
\eqref{eq:burnol-poly-rem} has weighted Fourier norm at
most $CR^{-M-1/2}e^{-5R/2}$; applying
\eqref{eq:burnol-cutoff-sampling} proves its second term.

\paragraph{The two-jet family and the all-window lower rate.}
For $P(t)=1-2t$, $M=0$ and $a=1$. Thus, at
$L=c\log\log R$ with any fixed $c>0$,
\[
 q((1-\chi_L)H_R)\sim
 m_\gamma A_{\chi,0}^2
             \frac{c^2(\log\log R)^2}{R},
 \qquad
 \|\chi_LH_R\|_2^2\sim
 B_{\chi,0}\frac{c\log\log R}{R}.
\]
The same asymptotics hold for the exact target-canceling
combination
\[
 Z_R=\mathcal JX_0^{e^{-R}}-c_R\mathcal JX_1^{e^{-R}},
 \qquad
 c_R=
 \frac{\widehat{\mathcal JX_0^{e^{-R}}}(\gamma)}
      {\widehat{\mathcal JX_1^{e^{-R}}}(\gamma)}
     =2+O_\gamma(R^{-1}).
\]
In fact $q(Z_R)=O_\gamma(R^{-1})$, and all $M=0$ cutoff
error estimates above are uniform over bounded polynomial
coefficients. Its local polynomial
$1-c_Rt$ has the same nonzero constant term and the same
leading cutoff samples and physical mass. No exact
integral constraint is needed for that cutoff calculation;
the uncut bound follows explicitly from
\[
 \|Z_R\|_q\le\|H_R\|_q+|c_R-2|\,
                   \|\mathcal JX_1^{e^{-R}}\|_q
                =O_\gamma(R^{-1/2}),
\]
using Proposition~\ref{prop:fd-burnol-weil}.

The proved all-window estimate of
Theorem~\ref{thm:cosine-lower-rate} is
\[
 a_1(L)\ge
 \exp[-C e^{2L/d_{\rm MT}}L(1+\log L)].
\]
At $L=c\log\log R$ its right side is
\[
 \exp\!\left[-C(\log R)^{2c/d_{\rm MT}}\,
        c\log\log R\bigl(1+\log(c\log\log R)\bigr)\right].
\]
The Rayleigh quotient of the compact two-jet
function is asymptotic to
\[
 m_\gamma\frac{A_{\chi,0}^2}{B_{\chi,0}}\,c\log\log R.
\]
The cutoff Weil norm is asymptotic to a positive constant times
$(\log\log R)/\sqrt R$, whereas the uncut Weil norm is
$O_\gamma(R^{-1/2})$.
These formulas quantify the missing term for this fixed
family and window scale. They provide no bound uniform in
the ordinate or as the number of jet orders grows.
The original all-height integer-deficit estimate remains
unproved.

